\scanpage{31}
\mspage{15}
go back up to \unsure{derived|various|divided} hypotheses arithmetico-geometric \inserted{in nature}, one can drop the anabelian aspects, to play to the full on the abelian aspects -- but in an arithmetic context. Instead of having \inserted{and using} only, as in corollary 1, a), \unsure{the hypothesis|of the hypothesis} that $H_1(\overline f)=H_1(\overline g);H_1(\overline X)\unsure{-|\longrightarrow}H^1(\overline Y\unsure{)}$ one has a finer hypothesis, on the equality of two homomorphisms of \uline{extensions} of $\mathcal E_{\overline K/K}$ by $H_{\unsure{1|n}}(\overline X,\widehat{\mathbb Z}_\ell)$ resp $H_{\unsure{1|n}}(\overline Y,\widehat{\mathbb Z}_\ell)$. $\forall\ell\in\mathbb P$.

Proof of corollary 3.
It suffices to see that $f(x)=g(x)$ for \unsure{$\forall$|every} $x\in X(K)$ For, after possibly \deleted{\unsure{extending}} \inserted{making a finite extension} of the base field, this remains true \deleted{then \unsure{$\forall$|every} extension of \deleted{\unsure{base}} \inserted{\unsure{field|field of}}} \inserted{for $x\in X(\overline K)$).} \deleted{\unsure{make|finite}}

Now a point defines a \inserted{conjugacy} class of \deleted{on the other hand} \unsure{Splittings|splittings} of $\mathcal E_{\overline X/K}$ over $\mathcal E_{\overline K/K}$, and a fortiori of the associated extension with abelian kernel -- the hypothesis implies that $f(x)$ and $g(x)$ define the $\unsure{=|\overline m}$ class of
\margin{30}
